\documentclass[../../main.tex]{subfiles}

\begin{document}
\chapter{Monte Carlo Methods}
\makeSubfileHeader{Monte Carlo Methods}

\section{Introduction}
    In this chapter, the Monte Carlo methods are discussed. Unlike DP, Monte Carlo does not
    require a perfect MDP model. Monte Carlo learns from experience, or simulated experience,
    if a model is ready.

    Monte Carlo methods are ways of solving reinforcement problem based on averaging sample returns.
    Unfortunately, the Monte Carlo method can be used only on episodic tasks.

\section{Monte Carlo Prediction}
    When a policy is given, we can estimate a state-value function of the policy by \textbf{Monte Carlo prediction}.
    Monte Carlo prediction averages the sample value of desired states and use it as the estimate of state-value function.
    The averaging is done when each episode terminates, unlike ordinary DP method which updates in each step.
    This idea of sampling and averaging the returns underlies all Monte Carlo methods.

    Precisely, Monte Carlo prediction are devided into two methods; \emph{first-visit MC method} and \emph{every-visit MC method}.
    The former one estimates $v_\pi(s)$ as the average of the returns first visits to $s$. The latter one averages the returns following all visits to $s$.
    These two MC methods are very similar but have slightly different theoretical properties.

    \Algorithm{firstVisitMCPredict}
    {First-Visit MC Prediction, for estimating $v_\pi$}
    {\inputitem{$\pi$}{a given policy}}
    {\inputitem{$V \simeq v_\pi$}{an estimate of $v_\pi$}}
    {
        \tcp{initialize}
        $V(s) \in \mathbb{R} \leftarrow \text{arbitrary values, for all } s \in \mathcal{S}$\;
        $R(s) \leftarrow \text{empty list, for all } s \in \mathcal{S}$\;

        \Repeat{convergence}
        {
            generate an episode following $\pi$: $S_0, A_0, R_1, S_1, A_1, \dots, S_{T-1}, A_{T-1}, R_T$\;
            $G \leftarrow 0$\;
            \ForEach{step of episode, $t = T-1, T-2, \dots, 1, 0$}
            {
                $G \leftarrow \gamma G + R_{t+1}$\;
                \tcp{in every-visit MC method, this if statement is ignored}
                \If{$S_t$ did not appear in $S_0, S_1, S_2, \dots, S_{t-1}$}
                {
                    append $G$ to $R(s)$\;
                    $V(S_t) \leftarrow average(R(s))$\;
                }
            }
        }
    }

    Due to the law of large numbers, the first-visit MC method and every-visit MC method converges to $v_\pi(s)$ as the number of visits to $s$ goes to infinity.

    Note that the MC method does not bootstrap; DP method uses values of other states when estimating a value of a given state. MC does not.
    Due to this property of MC method, MC method is useful when one requires only values of a subset of states.

\section{Monte Carlo Estimation of Action Values}
    If a model is not available, it is particularly useful to estimate action values rather than state values,
    since we do not know which action gives better return when a model is unavailable.
    The value we estimate now changes from $v_\pi(s)$ to $q_\pi(s,a)$. We just have to change state to state-action pair in algorithm \ref{alg:firstVisitMCPredict}.
    
    The only complication is that many state-action pairs won't be visited. If $\pi$ is a deterministic policy, we cannot get the
    values of action rather than $\pi(s)$. The problem of \emph{exploiting and exploring} rises here again. One way to maintain the exploration is to
    make the episode start with desired state-action pair, and every state-action pair has nonzero probability of being selected.
    We call this the assumption of \emph{exploring starts}. It is useful, but notice that it cannot be used in all situations; we cannot manually set the
    start when interacting with real-life environment.

\section{Monte Carlo Control}
    After the evaluation of policy, we can improve our policy, as we did in DP. The difference is that we have action-value function $q_\pi(s,a)$ instead of $v_\pi(s)$.
    This means that we do not need model to figure out the greedy policy.

    \subsection{Policy Improvement}
        For policy improvement, we select a greedy policy with respect to current action-value function $q_\pi(s,a)$.
        \begin{equation}
            \pi'(s) = \arg\max_a q_\pi(s,a)
        \end{equation}

        Note that this new policy $\pi'$ satisfies the conditions of policy improvement theorem;
        \begin{align}
            q_{\pi}(s, \pi'(s)) &= \max_a{q_{\pi}(s,a)} \\
            &\geq q_{\pi}(s, \pi(s)) \\
            &\geq v_\pi(s).
        \end{align}

    \subsection{Dealing with Infinite Iterations}

    \subsection{Algorithms}
    % add Monte Carlo ES algorithm
    % add On-policy first visit MC Control for \epsilon-soft polices algorithm


\end{document}